% ============================================================
\section{Iteración 3 --- DRV-05: motor de reglas}\label{sec:it03}
% ============================================================

% ------------------------------------------------------------
\subsection{Driver y alcance}

\begin{tabla}{|L{3.0cm}|Y|}
\fichacab
\parte{Driver}{DRV-05 --- Un único motor de reglas, parametrizable y ejecutable sin interfaz.}
\parte{Enunciado}{Existe una sola implementación del motor de reglas, con parámetros en configuración externa, separada de la interfaz, de la red y de la base de datos, determinista y capaz de reproducir partidas completas sin interfaz.}
\parte{Por qué le toca este turno}{Empata en 12 con DRV-03 y se ordena antes por dependencia: el reparto del segundo no se puede medir mientras el trabajo que hace el motor en cada jugada no esté acotado. Además es lo que el arbitraje de la iteración~1 consulta en cada jugada y lo que CD-11 necesita para entregar las jugadas legales.}
\parte{Qué decide esta iteración}{Cómo se construye el motor para que las reglas cambien sin tocar código y para que una partida completa se pueda ejecutar y reproducir sin abrir la interfaz.}
\parte{Decisiones ya tomadas}{DEC-01 a DEC-03 y lo decidido en las iteraciones~1 y~2. Pesa sobre todo DEC-02: con la lógica solo en el servidor, el motor tiene un único consumidor, así que el problema deja de ser que dos implementaciones coincidan y pasa a ser que una sola sea parametrizable y reproducible.}
\parte{Qué no decide}{Cuánto tiempo puede tardar el motor dentro del segundo, que se reparte en la iteración~4. Tampoco decide el contenido de las reglas, que es trabajo del diseñador (STK-06) y depende de Q-08, ni si habrá IA (Q-01).}
\end{tabla}

% ------------------------------------------------------------
\subsection{Paso 1. Revisión de entradas}

\subsubsection*{Objetivos de diseño}

\begin{tabla}{|L{1.5cm}|L{5.0cm}|Y|}
\hline
\filacab \cab{ID} & \cab{Objetivo} & \cab{Origen y qué exige aquí} \\ \hline
\endhead
OBJ-08 & El diseñador equilibra el juego sin depender de los programadores. &
CRN-13, con STK-06. El balanceo ocurre decenas de veces durante el desarrollo, no una vez. \\ \hline
OBJ-09 & El juego es original frente a Quoridor. &
CON-15 y Q-08, con STK-17 y STK-18. Si hay que introducir variantes propias, el motor tiene que admitirlas como parámetros y no como código nuevo. \\ \hline
OBJ-10 & Cada incremento semanal se puede probar sin jugar a mano. &
CRN-17 y CON-14, con STK-10 y STK-16. Sin ejecución sin interfaz no hay regresión posible en 14 semanas. \\ \hline
OBJ-11 & Un niño de 8 años entiende por qué una jugada no es válida. &
CRN-01 y CON-12, con STK-01 y STK-09. El motor no puede limitarse a decir que no. \\ \hline
\end{tabla}

\subsubsection*{Requisitos funcionales primarios}

\begin{tabla}{|L{1.3cm}|L{4.2cm}|Y|}
\hline
\filacab \cab{CU} & \cab{Caso de uso} & \cab{Qué exige del motor} \\ \hline
\endhead
CU-21 & Colocar un muro &
Las tres reglas más caras: que el muro no se solape ni se cruce (RN-04), que ningún jugador quede sin camino a su meta (RN-05) y que queden muros en la participación (RN-02). Además pide los surcos válidos y el efecto sobre el camino más corto (pasos 3 y 5). \\ \hline
CU-20 & Mover el peón &
Los desplazamientos legales, incluido el salto sobre el rival, y las casillas activables de la interfaz. \\ \hline
CU-25 & Finalizar la partida y aplicar la progresión &
Las condiciones de término (RN-01): meta alcanzada, reloj agotado, rendición, abandono o tablas. \\ \hline
CU-27 & Jugar contra la IA &
Cuatro niveles, tablero de 9$\times$9 o de 7$\times$7, reloj opcional, deshacer y sugerencias que explican la jugada recomendada. Todo ello son parámetros y consultas al mismo motor. \\ \hline
CU-37, CU-41 & Ver la repetición de una partida; avanzar en el tutorial &
Reproducen partidas jugada por jugada y posiciones preparadas, que solo funcionan si el motor es determinista. \\ \hline
\end{tabla}

\subsubsection*{Escenarios de atributos de calidad}

\begin{tabla}{|L{1.7cm}|L{3.3cm}|Y|}
\hline
\filacab \cab{Escenario} & \cab{Atributo} & \cab{Medida} \\ \hline
\endhead
ESC-12 \newline (ASR-06) & Mantenibilidad: modificabilidad &
El diseñador cambia el tablero y el número de muros en menos de 10 minutos, sin programador, sin recompilar y con 0 cambios de código. \\ \hline
ESC-17 & Mantenibilidad: capacidad de prueba &
500 partidas completas se ejecutan en menos de 10 minutos sin interfaz, el 100~\% de las repeticiones coincide con el original y cualquier fallo se reproduce señalando la jugada exacta. \\ \hline
ESC-02 & Usabilidad &
El motor tiene que poder decir qué camino quedaría cerrado, no solo que la jugada es inválida. Es una obligación sobre su interfaz. \\ \hline
ESC-01 \newline (ASR-01) & Desempeño &
Límite. El trabajo del motor por jugada entra en el segundo, y se mide en la iteración~4. \\ \hline
\end{tabla}

\subsubsection*{Restricciones}

\begin{tabla}{|L{1.9cm}|L{4.3cm}|Y|}
\hline
\filacab \cab{ID} & \cab{Restricción} & \cab{Cómo limita esta iteración} \\ \hline
\endhead
CON-15 & No infringir la propiedad intelectual de Quoridor & Puede obligar a variantes propias de las reglas con el proyecto avanzado. \\ \hline
CON-12 & El juego debe ser comprensible para niños de 8 años & El motivo del rechazo tiene que poder mostrarse de forma visual, lo que exige que el motor lo entregue. \\ \hline
CON-14 & Incrementos auditables cada semana & Exige regresión sin interfaz desde el primer incremento. \\ \hline
CON-01, CON-16 & C\# sobre .NET & El motor puede ser una biblioteca del mismo ensamblado que consume el servidor. \\ \hline
\end{tabla}

\subsubsection*{Decisiones de proyecto ya tomadas}

\begin{tabla}{|L{1.5cm}|L{4.3cm}|Y|}
\hline
\filacab \cab{ID} & \cab{Decisión} & \cab{Qué impone a esta iteración} \\ \hline
\endhead
DEC-02 & Toda la lógica del juego en el servidor & El motor tiene un solo consumidor. Desaparece el riesgo de dos implementaciones que discrepan y aparece la obligación de enumerar jugadas legales para el cliente. \\ \hline
CD-11 & El servidor entrega las jugadas legales junto con el turno & El motor no solo valida: tiene que enumerar y explicar. \\ \hline
CD-05 & Objeto activo por partida & El motor se ejecuta dentro del turno de una partida, sin estado compartido entre partidas. \\ \hline
\end{tabla}

\subsubsection*{Concerns}

\begin{tabla}{|L{1.5cm}|Y|}
\hline
\filacab \cab{ID} & \cab{Concern y qué aporta a la iteración} \\ \hline
\endhead
CRN-13 & \emph{Mientras equilibro el juego voy a cambiar el tamaño del tablero y el número de barreras decenas de veces. No puedo depender de que alguien recompile cada vez que quiero probar un valor.} \\ \hline
CRN-17 & El equipo de pruebas necesita ejecutar y repetir partidas enteras sin pasar por la interfaz. \\ \hline
CRN-12 & Cliente y servidor no pueden aplicar reglas distintas. Con DEC-02 el riesgo cambia de forma: ya no hay dos motores, pero sí dos vistas del mismo estado. \\ \hline
CRN-01 & El jugador quiere entender al momento por qué no puede hacer una jugada. \\ \hline
CRN-03 & El juego debe seguir interesando a quien ya sabe jugar, lo que se explora con variantes de reglas. \\ \hline
\end{tabla}

\subsubsection*{Preguntas abiertas que condicionan la iteración}

\begin{tabla}{|L{1.3cm}|Y|}
\hline
\filacab \cab{ID} & \cab{Cómo condiciona} \\ \hline
\endhead
Q-08 & Qué tan lejos deben estar las reglas de las de Quoridor. Si obliga a variantes, tienen que caber como parámetros. \\ \hline
Q-01 & Si habrá IA. CU-27 consulta al mismo motor para sugerir y para jugar, así que el motor debe poder evaluarse desde fuera de una partida. \\ \hline
\end{tabla}

% ------------------------------------------------------------
\subsection{Paso 2. Objetivo de la iteración y selección de entradas}

\subsubsection*{Priorización de las entradas}

\begin{tabla}{|L{3.2cm}|L{1.3cm}|Y|L{2.2cm}|}
\hline
\filacab \cab{Entrada} & \cab{Peso} & \cab{Por qué pesa lo que pesa} & \cab{Decisión} \\ \hline
\endhead
ESC-12 (ASR-06) & Alto & Es la medida del objetivo y decide si los parámetros viven en el código o fuera de él. & Se atiende \\ \hline
ESC-17 & Alto & Obliga a que el motor se pueda ejecutar sin interfaz y a que la partida sea determinista, que es una propiedad estructural, no un detalle de implementación. & Se atiende \\ \hline
CRN-13, CRN-17 & Alto & Son el origen de las dos medidas y describen la frecuencia real del cambio. & Se atiende \\ \hline
CU-20, CU-21 & Alto & Definen qué reglas tiene que aplicar y qué tiene que saber enumerar. & Se atiende \\ \hline
CD-11 & Alto & Obliga a que la interfaz del motor entregue jugadas legales y motivos, no solo un sí o un no. & Se atiende \\ \hline
ESC-02 & Medio & Aporta la obligación de explicar el rechazo; la parte visual es de la interfaz. & Se atiende como obligación \\ \hline
CU-27, CU-37, CU-41 & Medio & Consumen el motor fuera de una partida en línea y confirman que tiene que poder usarse por separado. & Se atiende \\ \hline
CON-15, Q-08 & Medio & Vuelven probable el cambio de reglas, que es lo que la parametrización protege. & Se atiende como supuesto \\ \hline
ESC-01 & Medio, como límite & El costo del motor por jugada entra en el segundo, pero su reparto se decide después. & Se difiere a la iteración~4 \\ \hline
CU-25 & Bajo & Las condiciones de término ya están decididas en la máquina de estados de CD-04; aquí solo se consultan. & Se atiende como supuesto \\ \hline
\end{tabla}

\subsubsection*{Objetivo de la iteración}

\begin{tabla}{|L{3.0cm}|Y|}
\fichacab
\parte{Objetivo}{Que las reglas del juego se puedan cambiar sin tocar código y que una partida completa se pueda ejecutar y reproducir sin abrir la interfaz.}
\parte{Driver que cubre}{DRV-05.}
\parte{Medida}{La de ESC-12: el diseñador cambia el tablero de 9$\times$9 a 7$\times$7 y los muros de 10 a 8 en menos de 10 minutos, sin programador, sin recompilar y con 0 cambios de código en el repositorio. Con ESC-17 como segunda comprobación de la misma propiedad: 500 partidas se ejecutan y se reproducen sin interfaz en menos de 10 minutos, con el 100~\% de coincidencia.}
\parte{Límite que respeta}{ESC-01. El trabajo del motor por jugada tiene que caber en el segundo, aunque su reparto se decida en la iteración~4.}
\parte{Incremento que produce}{El diseñador cambia los parámetros de una partida y juega con ellos sin que nadie recompile, y el equipo de pruebas lanza desde la línea de comandos las partidas registradas hasta ese momento y obtiene el mismo resultado que tuvieron. Las dos cosas se pueden demostrar en el incremento de la semana.}
\parte{Criterio de éxito}{La iteración termina cuando el cambio de parámetros se hace sin código y las 500 repeticiones coinciden.}
\end{tabla}

% ------------------------------------------------------------
\subsection{Paso 3. Elemento a descomponer}

% ------------------------------------------------------------
\Needspace*{0.4\textheight}
\subsubsection*{EL-03 --- Motor de reglas}

\begin{tabla}{|L{2.6cm}|Y|}
\fichacab
\parte{Elemento}{El motor de reglas: la pieza que, dada una posición y una jugada propuesta, dice si es legal, qué posición resulta y qué jugadas legales hay disponibles.}
\parte{Qué abarca}{Los desplazamientos del peón y el salto (CU-20), la colocación de muros con sus tres reglas (CU-21, RN-02, RN-04 y RN-05), el cálculo del camino más corto de cada jugador, la enumeración de jugadas legales, el motivo de cada rechazo y las condiciones de término que consulta CU-25. No abarca el turno, el reloj ni la sesión, que son del servicio de estado de la iteración~1.}
\parte{Por qué se elige}{Porque las dos medidas del objetivo dependen de él y de nada más: que las reglas cambien sin código depende de dónde vivan sus parámetros, y que una partida se reproduzca depende de que el motor no tenga azar ni reloj dentro. Se elige frente al servicio de estado de la partida, que ya tiene su iteración y que consume el motor sin definirlo, y frente a la interfaz del cliente, que con DEC-02 no contiene ninguna regla.}
\parte{Qué falta decidir sobre él}{Dónde viven el tamaño del tablero, el número de muros y el resto de valores que el diseñador cambia. Qué interfaz expone: si solo valida o además enumera y explica. Y qué se hace con todo lo que introduce variación, el azar del sorteo inicial o el reloj, que es lo que impide reproducir una partida.}
\parte{Evidencia en los casos de uso}{CU-21 describe en su paso 13 y su paso 14 las comprobaciones de solape y de camino, y en sus pasos 3 y 5 la necesidad de conocer los surcos válidos y el efecto sobre el camino más corto. CU-27 pide tablero de 9$\times$9 o de 7$\times$7 y sugerencias que expliquen la jugada recomendada, lo que solo es posible si el motor evalúa posiciones fuera de una partida. CU-37 reproduce una partida jugada por jugada. Ninguno dice dónde viven esos valores ni qué ocurre con el azar.}
\parte{Trazabilidad}{DRV-05. DEC-02. CD-05, CD-11. ESC-12 (ASR-06), ESC-17, ESC-02, con ESC-01 como límite. CU-20, CU-21, CU-25, CU-27, CU-37, CU-41. CON-15, CON-12, CON-14. CRN-13, CRN-17, CRN-12, CRN-01, CRN-03. Q-08, Q-01. STK-06, STK-10, STK-05, STK-09.}
\end{tabla}

% ------------------------------------------------------------
\subsection{Paso 4. Conceptos de diseño}

% ------------------------------------------------------------
\Needspace*{0.3\textheight}
\subsubsection*{CD-18 --- Motor como biblioteca sin dependencias de interfaz, red ni datos}

\begin{tabla}{|L{2.6cm}|Y|}
\fichacab
\parte{Estilo}{Estilo: separación de responsabilidades por dependencias.}
\parte{Descripción}{El motor es una biblioteca que no conoce la pantalla, el socket ni la base de datos: recibe una posición y devuelve un resultado. Atiende ESC-17, porque lo que no depende de la interfaz se puede ejecutar sin ella, y hace posible que CU-27, CU-37 y CU-41 lo usen fuera de una partida en línea. Se elige frente a un motor que forme parte del servicio de estado, que obligaría a levantar el servidor completo para probar una regla.}
\parte{Efectos}{La regresión de 500 partidas se lanza desde la línea de comandos. El servicio de estado consume el motor como una dependencia más. Obliga a definir con qué representación de la posición se habla con él, que es una decisión de interfaz que hay que tomar antes de construir.}
\parte{Trazabilidad}{DRV-05, DRV-08. CD-03. ESC-17, ESC-12. CU-20, CU-21, CU-27, CU-37, CU-41. CON-14, CON-01. CRN-17, CRN-12. STK-10, STK-05.}
\end{tabla}

% ------------------------------------------------------------
\Needspace*{0.3\textheight}
\subsubsection*{CD-19 --- Parámetros de reglas en configuración externa, fijados por partida}

\begin{tabla}{|L{2.6cm}|Y|}
\fichacab
\parte{Estilo}{Patrón de configuración.}
\parte{Descripción}{El tamaño del tablero, el número de muros por jugador y el resto de valores ajustables viven fuera del código, en un archivo de configuración, y el servidor los fija al crear cada partida. Atiende ESC-12, que es la medida del objetivo. Se elige frente a tenerlos como constantes, que obliga a recompilar y a un programador en cada prueba de balanceo (CRN-13), y frente a variantes codificadas una por una, que no cubren los valores que nadie previó.}
\parte{Efectos}{El diseñador cambia valores sin tocar el repositorio, y los cambios que Q-08 pueda imponer por CON-15 entran como parámetros. Cada partida guarda con qué parámetros se jugó, porque de otro modo su repetición no coincidiría. Obliga a validar los parámetros al cargarlos: un tablero imposible tiene que fallar al inicio y no a media partida.}
\parte{Trazabilidad}{DRV-05, DRV-08. ESC-12 (ASR-06), ESC-17. CU-21 (RN-02), CU-27, CU-37. CON-15, CON-13. CRN-13, CRN-03. Q-08. STK-06, STK-04.}
\end{tabla}

% ------------------------------------------------------------
\Needspace*{0.3\textheight}
\subsubsection*{CD-20 --- Motor determinista: sin reloj ni azar por dentro}

\begin{tabla}{|L{2.6cm}|Y|}
\fichacab
\parte{Estilo}{Patrón de diseño: función pura sobre el estado.}
\parte{Descripción}{El motor no consulta la hora ni genera números al azar: lo que necesite de fuera entra como dato, incluida la semilla del sorteo inicial. Atiende ESC-17, porque la misma secuencia de jugadas sobre la misma posición inicial da siempre el mismo resultado. Se elige frente a dejar que el motor consulte el reloj del sistema, que es lo que hace irreproducible una partida y convierte un fallo intermitente en algo que nadie puede investigar.}
\parte{Efectos}{El reloj de la partida vive en el servicio de estado (CD-04) y llega al motor como un valor. La repetición de CU-37 y la regresión de ESC-17 pasan a ser posibles. Obliga a registrar la semilla junto con la partida.}
\parte{Trazabilidad}{DRV-05, DRV-01. CD-04, CD-18. ESC-17, ESC-12. CU-20, CU-21, CU-37. CON-14. CRN-17. STK-10, STK-06.}
\end{tabla}

% ------------------------------------------------------------
\Needspace*{0.3\textheight}
\subsubsection*{CD-21 --- Interfaz que enumera las jugadas legales y explica el rechazo}

\begin{tabla}{|L{2.6cm}|Y|}
\fichacab
\parte{Estilo}{Patrón de interfaz de componente.}
\parte{Descripción}{El motor expone tres operaciones: validar una jugada, enumerar las jugadas legales de la posición y explicar por qué una jugada concreta no lo es, indicando qué camino quedaría cerrado o con qué muro se cruza. Atiende CD-11, que necesita la enumeración para el mensaje de turno, y ESC-02, que necesita el motivo para poder mostrarlo. Se elige frente a una interfaz que solo responda sí o no, que obligaría al cliente a deducir el resto y contradiría DEC-02.}
\parte{Efectos}{El mensaje de turno crece, con su costo para el protocolo y para el presupuesto de la iteración~4. La interfaz del cliente de CU-21 puede resaltar surcos válidos y mostrar el aviso de FA-02 sin calcular nada. Las sugerencias de CU-27 se apoyan en la misma enumeración.}
\parte{Trazabilidad}{DRV-05, DRV-01. DEC-02. CD-11. ESC-02, ESC-12, ESC-01. CU-20 (paso 1), CU-21 (pasos 3 y 5, FA-02), CU-27. CON-12. CRN-01, CRN-12. STK-01, STK-09, STK-06.}
\end{tabla}

% ------------------------------------------------------------
\Needspace*{0.3\textheight}
\subsubsection*{CD-22 --- Arnés de ejecución por lotes sin interfaz}

\begin{tabla}{|L{2.6cm}|Y|}
\fichacab
\parte{Estilo}{Táctica de capacidad de prueba: ejecución controlada del componente.}
\parte{Descripción}{Un programa aparte carga las partidas registradas, las vuelve a ejecutar jugada por jugada contra el motor y compara el resultado con el original, señalando la partida y la jugada donde aparece la primera diferencia. Atiende ESC-17 y lo convierte en la prueba del incremento semanal que pide CON-14. Se elige frente a probar a través de la interfaz, que es lo que CRN-17 declara inviable por tiempo.}
\parte{Efectos}{Cada incremento semanal se cierra con una ejecución de la regresión, que es la evidencia que STK-16 audita. Obliga a que las partidas queden registradas jugada por jugada, que es lo que la iteración~5 persistirá. Un fallo se reproduce sin abrir el cliente.}
\parte{Trazabilidad}{DRV-05, DRV-08. CD-18, CD-20. ESC-17, ESC-12. CU-37, CU-20, CU-21. CON-14, CON-13. CRN-17, CRN-08. STK-10, STK-16, STK-03.}
\end{tabla}
